% Figura conceptual editable. Las conexiones no representan magnitudes.
\begingroup
\setstretch{1}
\begin{tikzpicture}[
  >=Latex,
  every node/.style={inner sep=0pt,outer sep=0pt,align=left,text=introDark},
  problemaEtiqueta/.style={font=\fontfamily{phv}\selectfont\fontsize{7}{9}\selectfont,text=udaGray},
  problemaTitulo/.style={font=\fontfamily{phv}\selectfont\bfseries\fontsize{11}{13}\selectfont},
  problemaTexto/.style={font=\fontfamily{phv}\selectfont\fontsize{8.8}{11.5}\selectfont,text=udaGray},
  problemaNumero/.style={font=\fontfamily{phv}\selectfont\bfseries\fontsize{8}{10}\selectfont,text=introBlue}
]
  \path[use as bounding box] (0,0) rectangle (15,11);

  \node[anchor=north west,problemaEtiqueta] at (0.1,10.85) {OBJETIVOS ESPECÍFICOS};
  \draw[udaSky,line width=0.5pt] (0.1,10.42) -- (4.2,10.42);
  \node[anchor=north west,font=\fontfamily{phv}\selectfont\fontsize{17}{20}\selectfont] at (0.1,10.05)
    {De los indicadores\\\textbf{a la validación comparativa}};

  % Semicírculo y arco de conexión: síntesis cualitativa de las condiciones.
  \fill[introBlue!12!white] (8.35,8.30)
    arc[start angle=90,end angle=270,radius=3.35] -- cycle;
  \draw[udaSky,line width=0.7pt] (8.35,8.60)
    arc[start angle=90,end angle=270,radius=3.65];
  \draw[udaSky,line width=0.6pt] (4.15,8.60) -- (8.35,8.60);
  \draw[udaSky,line width=0.6pt] (4.40,4.95) -- (4.70,4.95);
  \draw[udaSky,line width=0.6pt] (4.15,1.30) -- (8.35,1.30);
  \foreach \x/\y in {8.35/8.60,4.70/4.95,8.35/1.30}{
    \fill[introBlue] (\x,\y) circle (0.065);
  }
  \draw[udaSky,line width=0.5pt] (8.35,8.30) -- (8.35,1.60);

  % Condiciones del caso de estudio. Los números solo ordenan la lectura.
  \node[anchor=north west,problemaNumero] at (0.10,8.35) {01};
  \node[anchor=north west,problemaTitulo] at (0.62,8.35) {Configurar la matriz};
  \node[anchor=north west,problemaTexto] at (0.62,7.73)
    {Integrar cinco marcos\\y definir indicadores.};

  \node[anchor=north west,problemaNumero] at (0.10,5.65) {02};
  \node[anchor=north west,problemaTitulo] at (0.62,5.65) {Evaluar el caso base};
  \node[anchor=north west,problemaTexto] at (0.62,5.03)
    {Aplicar la matriz\\y reconocer las brechas\\del proyecto.};

  \node[anchor=north west,problemaNumero] at (0.10,2.58) {03};
  \node[anchor=north west,problemaTitulo] at (0.62,2.58) {Desarrollar la propuesta};
  \node[anchor=north west,problemaTexto] at (0.62,1.96)
    {Integrar estrategias pasivas\\y materiales sostenibles.};

  % Problema central, situado sobre la superficie semicircular.
  \node[anchor=north west,problemaEtiqueta,text=introBlue] at (5.57,6.22) {CUARTO OBJETIVO};
  \node[anchor=north west,font=\fontfamily{phv}\selectfont\bfseries\fontsize{11.5}{14}\selectfont] at (5.57,5.72)
    {Validar la\\incidencia de\\la propuesta};
  \node[anchor=north west,problemaTexto] at (5.57,3.86)
    {Comparación de\\energía y confort.};
  \draw[->,introOrange,line width=1pt] (8.48,5.45) -- (9.03,5.45);

  % Evaluación requerida: mismo alcance del texto de la investigación.
  \node[anchor=north west,problemaEtiqueta,text=introBlue] at (9.20,8.16) {PRODUCTOS VERIFICABLES};
  \node[anchor=north west,font=\fontfamily{phv}\selectfont\bfseries\fontsize{16}{19}\selectfont] at (9.20,7.60)
    {Conectar método\\y resultados};
  \node[anchor=north west,problemaTexto] at (9.20,6.12)
    {Conservar la trazabilidad\\entre objetivos y evidencias.};

  \foreach \y in {4.67,3.15,1.63}{
    \path[draw=udaSky!75!white,fill=introSoft,line width=0.55pt,rounded corners=5pt]
      (9.20,\y-0.60) rectangle (14.85,\y+0.60);
    \fill[white] (9.77,\y) circle (0.29);
  }
  \node[anchor=west,problemaTitulo] at (10.31,4.87) {Matriz y diagnóstico};
  \node[anchor=west,problemaTexto] at (10.31,4.43) {Criterios y línea base};
  \node[anchor=west,problemaTitulo] at (10.31,3.35) {Propuesta de mejora};
  \node[anchor=west,problemaTexto] at (10.31,2.91) {Escenarios documentados};
  \node[anchor=west,problemaTitulo] at (10.31,1.83) {Validación comparativa};
  \node[anchor=west,problemaTexto] at (10.31,1.39) {Resultados y limitaciones};

  % Iconos lineales de energía, temperatura y materialidad.
  \draw[introBlue,line width=0.85pt,line join=round]
    (9.81,4.87) -- (9.64,4.65) -- (9.77,4.65) -- (9.72,4.46) -- (9.92,4.71) -- (9.79,4.71) -- cycle;
  \draw[introBlue,line width=0.8pt,rounded corners=1.2pt]
    (9.71,3.11) -- (9.71,3.34) -- (9.82,3.34) -- (9.82,3.11);
  \draw[introBlue,line width=0.8pt] (9.765,3.04) circle (0.105);
  \draw[introBlue,line width=0.75pt] (9.765,3.10) -- (9.765,3.25);
  \draw[introBlue,line width=0.8pt,line join=round]
    (9.56,1.70) -- (9.77,1.83) -- (9.98,1.70) -- (9.77,1.57) -- cycle;
  \draw[introBlue,line width=0.8pt,line join=round]
    (9.56,1.60) -- (9.77,1.47) -- (9.98,1.60);

  \node[anchor=south west,problemaEtiqueta] at (0.10,0.15)
    {DEFENSORÍA DEL PUEBLO · SAN CRISTÓBAL};
  \node[anchor=south east,problemaEtiqueta] at (14.85,0.15)
    {Configurar \quad / \quad Evaluar \quad / \quad Desarrollar \quad / \quad Validar};
\end{tikzpicture}
\endgroup
